\section{Demonstration}
In this section, the demonstration phase of the DSRM methodology is discussed. It illustrates how the proposed framework can be applied to address the research problem identified in Section 5 within a realistic software development context.

The research proposal includes the development of a process-oriented framework inspired by Agile and Scrum principles, aimed at providing actionable guidance for organizations seeking to adopt AI and LLM-based tools in a secure, responsible, sustainable, and value-driven manner. To demonstrate how the framework can be used in practice, a realistic project scenario is considered based on an next-generation consultancy.

This consultancy is a next-generation consultancy that helps organizations solve business problems with technology rapidly, creatively, and always with people at the centre. Unlike traditional consultancies or software houses, this consultancy prioritizes business value, high-speed delivery, and new ways of building, including Vibe Coding and AI-assisted prototyping. Their mission is to unlock progress where other approaches stall.

In the demonstration scenario, a team is tasked with developing a small to medium-sized software application. The organization already pioneers the use of AI, LLMs, and other AI-augmented tools to accelerate software development through practices such as Vibe Coding. The team operates using Scrum-based practices, short feedback cycles, and continuous alignment with stakeholders. However, despite this advanced adoption, this consultancy, like many modern consultancies, continues to seek new methodological evolutions of Scrum and Agile that can systematically standardize and govern the use of AI across the entire Software Development Lifecycle. In the absence of clear guidance and structured processes for integrating AI into development workflows, the team encounters challenges that are typical of AI related environments:
\begin{enumerate}
    \item difficulty validating and reviewing AI-generated code in a consistent and structured manner;
    \item increased pressure on requirements refinement, as rapid prototyping exposes ambiguities earlier in the process;
    \item increased complexity in functional and security testing due to unpredictable AI behaviour;
    \item unclear responsibility boundaries between human contributors and AI-generated artifacts.
\end{enumerate}

These challenges reflect the gaps identified in the systematic literature review and motivate the application of the proposed framework. Within this scenario, the framework is applied as the primary process structure to guide the adoption and integration of AI tools, define explicit human–AI collaboration practices, introduce human-in-the-loop quality assurance checkpoints, and guide decision-making across SDLC phases. Development artifacts such as AI-generated outputs, prompt formulations, review notes, and iteration-level observations are used to illustrate how the framework governs AI participation and supports accountable development practices.

The demonstration will be structured around a limited number of short development iterations aligned with Agile cycles. Across these iterations, the application of the framework focuses on observing decision points related to AI tool selection, prompt formulation, validation practices, and responsibility allocation between human developers and AI systems. The demonstration emphasizes process behavior rather than delivery performance, enabling the analysis of how the framework shapes development practices, oversight mechanisms, and collaboration patterns.

The next section focuses on evaluating the framework’s suitability, usefulness, and alignment with the needs observed in this scenario, in accordance with the Evaluation stage of the DSRM methodology.
